% s03-1-1-basics.tex —— 3.1.1 小节：因子化基础（印刷页 56–60）

\chapter{Iwahori 基本局部等价}
\label{chap:iwahori}

\section{因子化}
\label{sec:factorization}

\subsection{因子化基础}
\label{sec:fact-basics}

本小节简要回顾因子化的一些定义与概念。目的主要是固定记号；详细讨论见 \cite{Ras}。

全文始终固定：$k$ 为特征 $0$ 的代数闭域，$X$ 为 $k$ 上光滑、射影、连通的曲线，
并固定一个闭点 $x_0\in X$。

先固定因子化对象所栖居的空间的记号。粗略地说，模式如下：我们的「代数型」对象
（因子化范畴、因子化代数、因子化空间）将栖居于 $\Ran$ 之上；而固定一个「代数对象」后，
相应类型的「模型型」对象（因子化模范畴、因子化模、因子化模空间）将栖居于
$\Ran_{x_0}$ 之上。这是因为在此文中，因子化「模型型」对象的预期应用是研究固定点 $x_0$
处的层级结构。

\begin{definition}
\label{def:ran}
记 $\Ran$ 为 $X$ 的\emph{带单位的} $\Ran$（范畴）预叠。对试验仿射概形 $S$，
\[
\Ran(S)=\{\text{$X_{dR}(S)$ 的有限子集}\},
\]
并把它看作按包含关系构成的偏序集。
\end{definition}
\srcpage{56}

\begin{remark}
\label{rem:ran-convention}
我们采用这样的约定：不加修饰的 $\Ran$ 指 $\Ran$ 空间的带单位版本；而其他文献可能
用同一记号指非带单位版本，其取值是 $X_{dR}(S)$ 有限子集构成的\emph{集合}而非偏序集
（并且可能排除空子集）。
\end{remark}

\yizhu{可对照 nLab 的「Ran space」词条：该词条把 Ran 空间定义为「全体非空有限子集」构成的空间，
与本文含空集的带单位版本不同。可见这不是本文独有的约定分歧，而是文献间的普遍现象。}{B：nLab, 「Ran space」}

今后凡与因子化有关的内容，除非明确说明，默认指带单位版本。一个例外是：因子化函子
（即因子化范畴与因子化模范畴之间的态射）允许是宽松带单位的（lax-unital），
而不必严格带单位。

上述预叠有一个变体 $\Ran^{\subset}$，它分类 $X$ 的有限子集的有序对，其中第一个含于
第二个。有投影映射
\[
\mathrm{pr}_{\mathrm{small}},\ \mathrm{pr}_{\mathrm{big}}\colon
\Ran^{\subset}\longrightarrow \Ran,
\]
分别把有序对 $A\subset B\subset X_{dR}(S)$ 送到 $A$ 与 $B$。

\begin{definition}
\label{def:ran-Y}
设 $Y\to \Ran$ 为态射。记
\[
\Ran_Y=Y\times_{\Ran,\mathrm{pr}_{\mathrm{small}}}\Ran^{\subset},
\]
其中纤维积沿 $\mathrm{pr}_{\mathrm{small}}$ 构造。于是有由
$\mathrm{pr}_{\mathrm{big}}\colon \Ran^{\subset}\to \Ran$ 诱导的映射
\[
\mathrm{pr}_{\mathrm{big},Y}\colon \Ran_Y\longrightarrow \Ran .
\]
\end{definition}
\srcpage{57}

对我们最重要的情形是 $Y=\pt$、$Y\to \Ran$ 为含 $x_0$ 的单点子集所给。此时记相应预叠为
$\Ran_{x_0}$。

\begin{remark}
\label{rem:ran-x0}
对试验仿射概形 $S$，$\Ran_{x_0}(S)$ 分类 $X_{dR}(S)$ 中含常值映射
$S\to \pt\xrightarrow{x_0}X_{dR}$ 的有限子集。因此
$\mathrm{pr}_{\mathrm{big},x_0}\colon \Ran_{x_0}\to \Ran$ 是嵌入。
\end{remark}

下面定义因子化范畴的概念。直观地说，这给出一种按族给每一点组 $x\subset X$ 指派范畴的方式，
且关心点之间的碰撞。考虑子集
\[
[\Ran\times \Ran]_{\mathrm{disj}}\subset \Ran\times \Ran
\]
分类 $X$ 的有限子集的互不相交对。再作图表
\[
\Ran\times \Ran\supset [\Ran\times \Ran]_{\mathrm{disj}}
\xrightarrow{\ \cup\ }\Ran .
\]

\begin{definition}
\label{def:fact-cat}
\emph{因子化范畴}是 $\Ran$ 上的范畴层 $\mathcal{C}$，连同因子化数据，即一个等价
\[
(\mathcal{C}\boxtimes \mathcal{C})|_{\mathrm{disj}}
\simeq \cup^{*}\mathcal{C},
\]
以及高阶同伦相容性。
\end{definition}

记因子化范畴构成的范畴为 $\FactAlgCat$。

为定义（$x_0$ 处的）因子化模范畴，考虑子集
\[
[\Ran\times \Ran_{x_0}]_{\mathrm{disj}}\subset \Ran\times \Ran_{x_0},
\]
其定义方式显然。同样可作图表
\[
\Ran\times\Ran_{x_0}\supset[\Ran\times\Ran_{x_0}]_{\mathrm{disj}}
\xrightarrow{\ \cup\ }\Ran_{x_0}.
\]
\srcpage{58}

\begin{definition}
\label{def:fact-modcat}
设 $\mathcal{C}$ 为因子化范畴。（$x_0$ 处的）\emph{因子化 $\mathcal{C}$-模范畴}是
$\Ran_{x_0}$ 上的范畴层 $\mathcal{M}$，连同 $\mathcal{C}$ 上的因子化模数据，即一个等价
\[
(\mathcal{C}\boxtimes \mathcal{M})|_{\mathrm{disj}}
\simeq \cup^{*}\mathcal{M},
\]
以及高阶同伦相容性。
\end{definition}

记因子化 $\mathcal{C}$-模范畴构成的范畴为
$\mathcal{C}\text{-}\FactModCat_{x_0}$。

\begin{remark}
\label{rem:pairs}
有时我们不想固定因子化范畴 $\mathcal{C}$，而要考虑对
$(\mathcal{C},\mathcal{M})$，其中 $\mathcal{M}$ 是因子化 $\mathcal{C}$-模范畴。
这种对构成的范畴记作 $\FactModCat_{x_0}$，其中的态射
\[
(\mathcal{C}_1,\mathcal{M}_1)\to(\mathcal{C}_2,\mathcal{M}_2)
\]
由因子化函子（即 $\FactAlgCat$ 中的态射）$\mathcal{C}_1\to \mathcal{C}_2$ 以及
（$\Ran_{x_0}$ 上）范畴层的态射 $\mathcal{M}_1\to \mathcal{M}_2$ 构成，
并与上述因子化函子相容。
\end{remark}

下面讨论对象层面的因子化结构。

\begin{definition}
\label{def:fact-alg}
设 $\mathcal{C}$ 为因子化范畴。$\mathcal{C}$ 中的\emph{因子化代数} $\underline{a}$
是 $\mathcal{C}$ 的整体截面，连同因子化数据，即一个同构
\[
(\underline{a}\boxtimes \underline{a})|_{\mathrm{disj}}
\simeq \cup^{!}\underline{a},
\]
以及高阶同伦相容性。
\end{definition}
\srcpage{59}

记 $\mathcal{C}$ 中因子化代数构成的范畴为 $\FactAlg(\mathcal{C})$。

\begin{definition}
\label{def:fact-mod}
设 $Y\to \Ran$。再设 $(\mathcal{C},\mathcal{M})\in \FactModCat_Y$
（其定义只需把 $\FactModCat_{x_0}$ 定义中的 $\Ran_{x_0}$ 换成 $\Ran_Y$）为对。
则对 $\mathcal{C}$ 中的因子化代数 $\underline{a}$ 及任意 $Z\to \Ran_Y$，
$Z$ 处 $\mathcal{M}$ 中的\emph{因子化 $\underline{a}$-模} $\underline{m}$
是 $\mathcal{M}$ 在 $Z$ 上的截面，连同 $\underline{a}$ 上的因子化模数据，即一个同构
\[
(\underline{a}\boxtimes \underline{m})|_{\mathrm{disj}}
\simeq \cup^{!}\underline{m},
\]
以及高阶同伦相容性。
\end{definition}

记此类对象构成的范畴为
$\underline{a}\text{-}\FactMod_{Z}(\mathcal{M})$。

我们之所以不只对 $x_0$ 而要对更一般的 $Z$ 定义因子化模，是因为想让 $Z$ 变动，
从而构造因子化模的因子化范畴与因子化模范畴。

\begin{definition}
\label{def:fact-mod-sheaf}
设 $(\mathcal{C},\mathcal{M})\in \FactModCat_{x_0}$，$\underline{a}\in \FactAlg(\mathcal{C})$。
考虑 $\Ran$ 上的范畴层
\[
\underline{a}\text{-}\FactMod(\mathcal{C}),
\]
它把 $\Ran$ 上的试验仿射概形 $S$ 送到范畴
$\underline{a}\text{-}\FactMod_{S}(\mathcal{C})$。它具有自然的因子化范畴结构。
再考虑 $\Ran_{x_0}$ 上的范畴层
\[
\underline{a}\text{-}\FactMod(\mathcal{M}),
\]
它把 $\Ran_{x_0}$ 上的试验仿射概形 $S$ 送到范畴
$\underline{a}\text{-}\FactMod_{S}(\mathcal{M})$。它具有自然的、$x_0$ 处
$\underline{a}\text{-}\FactMod(\mathcal{C})$-因子化模范畴结构。
\end{definition}

以上给出了构造因子化范畴与因子化模范畴的代数方法。本小节最后给出几何方法。
\srcpage{60}

\begin{definition}
\label{def:fact-space}
\emph{因子化空间}是空间 $Y\to \Ran$，连同因子化数据，其中包括一个同构
\[
(Y\times Y)\times_{\Ran\times \Ran}[\Ran\times \Ran]_{\mathrm{disj}}
\simeq Y\times_{\Ran,\cup}[\Ran\times \Ran]_{\mathrm{disj}}.
\]
\end{definition}

\begin{definition}
\label{def:fact-mod-space}
（$x_0$ 处的）\emph{因子化 $Y$-模空间}是空间 $Z\to \Ran_{x_0}$，
连同 $Y$ 上的因子化模数据，其中包括一个同构
\[
(Y\times Z)\times_{\Ran\times\Ran_{x_0}}[\Ran\times\Ran_{x_0}]_{\mathrm{disj}}
\simeq Z\times_{\Ran_{x_0},\cup}[\Ran\times\Ran_{x_0}]_{\mathrm{disj}}.
\]
\end{definition}

给定一个具有拉回与 Künneth 公式的「良好」层论（例如 $\QCoh$），因子化空间 $Y$
通过
\[
(S\to \Ran)\mapsto \QCoh(S\times_{\Ran}Y)
\]
给出因子化范畴；因子化 $Y$-模空间 $Z$ 通过
\[
(S\to \Ran_{x_0})\mapsto \QCoh(S\times_{\Ran_{x_0}}Z)
\]
给出因子化 $\QCoh(Y)$-模范畴。
